\chapter{Vowels Conditioned by Spelling}\label{ch:vowels}

\section{Observation}
A purely phonemic mapping from the pronouncing dictionary reproduces the page's consonants well and its vowels less well. Comparing the computed forms with the attested forms shows a regular departure: the letter used for a vowel sometimes follows the English spelling rather than the sound.

\begin{enumerate}
\item \textbf{Stressed o.} The dictionary gives \emph{Rollins}, \emph{on}, \emph{psychology} and \emph{concept} an open back vowel. The specimens write them with damma and waw: \ar{رُولِنْزْ}, \ar{أُونْ}, \ar{سَايْكُولُوجِي}, \ar{كُونْسِبْتْ}. Contrast \emph{part}, written \ar{بَارْتْ} with fatha and alef, where the spelling letter is $a$.
\item \textbf{Open-syllable e and i.} \emph{Helen} and \emph{psychocinema} are written with kasra and \ar{ي}: \ar{هِيلِنْ}, \ar{سِينِمَا}. The stressed vowel letter is followed by one consonant and another vowel. In \emph{delves}, \emph{this} and \emph{its}, where the vowel letter is followed by a cluster, only kasra is written.
\item \textbf{Schwa colour.} An unstressed vowel takes the colour of its letter: $a$ gives fatha, $e$ and $i$ give kasra, $o$ gives damma, as in \ar{سِمْبْتُمْزْ} \emph{symptoms}.
\item \textbf{Short initial vowels.} An unstressed vowel in the first syllable is written short, with no vowel letter: \ar{رِبْرِشَنْ} \emph{repression}, \ar{يُتُوبِيَانِيزْمْ} \emph{utopianism}.
\item \textbf{Hosted hiatus.} A vowel after \ar{ي} is written on the \ar{ي}, as in \ar{بِيَ} and \ar{ثِيُ}.
\end{enumerate}

\section{Why pronunciation and spelling both contribute}
The notation could have been defined as a pure phonetic transcription, with each vowel written from its dictionary sound alone. The author's practice shows that it is not. Familiar English words carry a visual identity that readers use in recognising them, and the notation keeps some of it. The word \emph{book} is written \ar{بُوكْ} with a waw because English writes two vowel letters, and the waw retains that feature even though it does not reproduce the short vowel narrowly. The waw is a deliberate orthographic choice. It is not an error of pronunciation to be corrected.

The consequence for later revision is practical. A departure from dictionary pronunciation in an attested form is not automatically a mistake. It may be the spelling showing through, and a rule that removes it would make the notation more phonetic and less like the author's writing. Before such a form is changed, the question to ask is whether the departure is systematic across words with the same spelling, as the waw of \emph{book}, \emph{good} and \emph{foot} is, or whether it is a lapse.

\section{Reading of the pattern}
These observations suggest that the notation is a hybrid. It is phonetic in its consonants and in the broad shape of its syllables, and orthographic in the choice of some vowel letters. The evidence is one author's page, a typed passage and generated texts, so the hybrid reading is a hypothesis about how the notation behaves, not a finding about how it must behave. The rules implemented from it are listed with their results in Chapter~\ref{ch:eval}.

\section{Remaining disagreements}
After the rules are applied, \NDiffers{} of the \NAttested{} attested forms still differ from the computed forms by more than marks. Table~\ref{tab:differs} lists them. They fall into five groups:
\begin{description}
\item[Lexical alef.] \emph{capitalist}, \emph{capitalism}, \emph{utopianism} and \emph{psychoanalytic} carry an alef where the rules give only fatha. Other words with the same spelling pattern, such as \emph{fundamental} and \emph{cultural}, do not. No rule from spelling or stress separates them, so these stay in the attested list.
\item[The choice of v.] \emph{individual} and \emph{universal}, written with \ar{ب} in the typed passage, discussed in Chapter~\ref{ch:evidence}. The default \ar{ف} is the author's decision and the specimens disagree with it in these two words.
\item[Failed alignment.] When the number of vowel letter groups differs from the number of vowel phonemes, as in \emph{ideology} and \emph{individual}, spelling-based colouring is skipped and the phonemic default applies.
\item[Seat letter.] The article \emph{a} is written \ar{اَ} in the typed passage and \ar{أَ} in the original fixed table. Both appear in the sources; the toolkit keeps the hamza form.
\item[Generated-text noise.] \emph{nature} in the typed passage, which uses \ar{ش} where the page has \ar{تْش}, and \emph{structures}, whose typed spelling departs from both the page conventions and the sound.
\end{description}

\begin{table}[ht]\centering\small
\begin{tabular}{@{}lll@{}}
\toprule
English & Attested & Computed\\
\midrule
a & \ar{اَ} & \ar{أَ}\\
capitalism & \ar{كَابِتَالِيزْمْ} & \ar{كَابِتَلِزْمْ}\\
capitalist & \ar{كَابِتَالِسْتْ} & \ar{كَابِتَلِسْتْ}\\
ideology & \ar{آيدِيُولُوجِي} & \ar{آيْدِيَالَجِي}\\
individual & \ar{إِنْدِبِدْيُوَالْ} & \ar{إِنْدَفِجَوَلْ}\\
nature & \ar{نَيْشَرْ} & \ar{نَيْتْشَرْ}\\
psychoanalytic & \ar{سَايْكُوأَنَالِيتِيكْ} & \ar{سَايْكُوأَنَلِتِكْ}\\
structures & \ar{سْتْرَكْشُورْزْ} & \ar{سْتْرَكْتْشَرْزْ}\\
universal & \ar{يُونِبِرْسَلْ} & \ar{يُونِفِرْسَلْ}\\
utopianism & \ar{يُتُوبِيَانِيزْمْ} & \ar{يُتُوبِيَنِزْمْ}\\
\bottomrule
\end{tabular}
\caption{Attested forms that differ from the computed forms in their letters.}\label{tab:differs}
\end{table}

\section{The rules in order}
The vowel rules are applied in a fixed order for each vowel phoneme, and the first rule that fires decides the mark and the vowel letters.
\begin{enumerate}
\item Schwa (AH) takes the colour of its spelling letter: $e$, $i$ and $y$ give kasra, $o$ gives damma, $a$ gives fatha, with fatha also for the vowel pairs \emph{ia} and \emph{io} and for word-final \emph{-em}. A stressed AH spelled with $o$ is written with damma and waw.
\item A rhotic vowel (ER) is written with fatha when unstressed, with the exception of the prefix \emph{inter-}, which keeps kasra. Stressed ER takes its mark from the spelling letter, but \emph{ea} gives fatha (\emph{learn}).
\item The vowels AA and AO, when the spelling letter is $o$, are written with damma and waw (\emph{on}, \emph{Rollins}).
\item A stressed AE that is not word-initial gains an alef (\emph{lack}).
\item A stressed EH or IH in an open syllable spelled with $e$ or $i$ gains a yeh (\emph{Helen}, \emph{cinema}).
\item An unstressed IY or UW in the first syllable is written short.
\item The vowel UH spelled with a written \emph{oo} or \emph{ou} keeps a waw (\emph{book}, \emph{good}, \emph{foot}).
\end{enumerate}
Each rule can be traced to specimen words, and each has exceptions among the attested forms that are listed in Table~\ref{tab:differs}.

\section{Double o and the word book}
The last rule came from a correction by the author. A composed example wrote \emph{book} as \ar{بُكْ}, following the sound, and the author preferred \ar{بُوكْ} on the ground that English writes two vowel letters and the notation should show them. The author also noted that without the waw an Arabic reader might hear a different vowel. The same reasoning applies to \emph{look}, \emph{good}, \emph{foot} and \emph{wood}, and the toolkit now writes all of them with a waw. It does not yet extend to every doubled vowel, because the specimens contain no attested case of other doubled letters, apart from \emph{ee}, which already receives a yeh as a long vowel.

\section{Hosted vowels after a glide}
When a vowel follows a glide, the mark is placed on the glide, so no separate hamza seat is needed. This accounts for \emph{flower}, \emph{power} and \emph{hour}, where the diphthong in the first syllable ends in a waw glide and the unstressed syllable takes its fatha or kasra on that waw. The same device was already used after \ar{ي}, as in \emph{utopianism}.

\section{Ng and the hard g}
In the dictionary, \emph{finger}, \emph{anger} and \emph{English} have an ng sound followed by a hard g. Written as separate sounds this gives a doubled ghain, which no specimen shows. The toolkit therefore writes the nasal and one ghain, as in \ar{إِنْغْلِشْ}. Words such as \emph{singer} and \emph{longer}, where the dictionary has only the nasal, are written the same way, which is a merger the notation accepts.

\section{Examples}
Table~\ref{tab:whybrid} gives words from the specimens and words that were not in them, written by the rules. The specimen words reproduce their attested forms. The others show how the rules extend.

\begin{table}[ht]\centering\small
\begin{tabular}{@{}lll@{}}
\toprule
English & Arabic script & Respelling\\
\midrule
psychology & \ar{سَايْكُولُوجِي} & \textit{saykaalujee}\\
concept & \ar{كُونْسِبْتْ} & \textit{kaansept}\\
Rollins & \ar{رُولِنْزْ} & \textit{raalinz}\\
on & \ar{أُونْ} & \textit{aan}\\
cinema & \ar{سِينِمَا} & \textit{sinumu}\\
Helen & \ar{هِيلِنْ} & \textit{helun}\\
repression & \ar{رِبْرِشَنْ} & \textit{reepreshun}\\
social & \ar{سُوشَلْ} & \textit{sohshul}\\
nation & \ar{نَيْشَنْ} & \textit{neyshun}\\
system & \ar{سِسْتَمْ} & \textit{sistum}\\
lack & \ar{لَاكْ} & \textit{lak}\\
part & \ar{بَارْتْ} & \textit{paart}\\
flower & \ar{فْلَاوَرْ} & \textit{flower}\\
hour & \ar{آوَرْ} & \textit{ower}\\
power & \ar{بَاوَرْ} & \textit{power}\\
tower & \ar{تَاوَرْ} & \textit{tower}\\
finger & \ar{فِنْغَرْ} & \textit{fingger}\\
English & \ar{إِنْغْلِشْ} & \textit{ingglish}\\
\bottomrule
\end{tabular}
\caption{Words that exercise the spelling-conditioned rules.}\label{tab:whybrid}
\end{table}
